%------------------------------------------------------------------------------%
\documentclass[fleqn,utf8,aspectratio=169,12pt]{beamer}

\usepackage{gaal}

\title{Dependência e Independência Linear}

%------------------------------------------------------------------------------%
\begin{document}

\addtitle

%------------------------------------------------------------------------------%
\section{Dependência e Independência Linear}

%------------------------------------------------------------------------------%
\begin{frame}{Combinação Linear}
  Sejam
  \[
    \vec{v}_1, \vec{v}_2, \ldots, \vec{v}_k \in \R^n
  \]
  \[
    \pause
    \alpha_1, \alpha_2, \ldots, \alpha_k \in \R
  \]

  \pause
  Uma \emph{combinação linear} é
  \[
    \pause
    \alpha_1\vec{v}_1+
    \alpha_2\vec{v}_2+
    \cdots+
    \alpha_k\vec{v}_k
  \]
\end{frame}

%------------------------------------------------------------------------------%
\begin{frame}{Dependência Linear}
  Os vetores
  \[
    \pause
    \vec{v}_1, \ldots, \vec{v}_k
  \]
  \pause
  são \emph{linearmente dependentes} se existem escalares
  \[
    \pause
    \alpha_1, \ldots, \alpha_k
  \]
  \pause
  \emph{não todos nulos} tais que
  \[
    \pause
    \alpha_1\vec{v}_1 + \cdots + \alpha_k\vec{v}_k
    \pause
    = \vec{0}
  \]
\end{frame}

%------------------------------------------------------------------------------%
\begin{frame}{Dependência Linear}
  Se
  \[
    \alpha_1\vec{v}_1 + \cdots + \alpha_k\vec{v}_k = \vec{0}
  \]
  \pause
  com algum \(\alpha_i \neq 0\)

  \pause
  \bigskip
  Podemos isolar o vetor \(\vec{v}_i\)
  \[
    \pause
    \vec{v}_i = \frac{\alpha_1}{\alpha_i}\vec{v}_1 + \cdots + \frac{\alpha_k}{\alpha_i}\vec{v}_k
  \]
  \pause
  e escreve-lo como combinação dos demais
\end{frame}

%------------------------------------------------------------------------------%
\begin{frame}{Independência Linear}
  Os vetores
  \[
    \pause
    \vec{v}_1, \ldots, \vec{v}_k
  \]
  \pause
  são \emph{linearmente independentes} se
  \[
    \pause
    \alpha_1\vec{v}_1 + \cdots + \alpha_k\vec{v}_k = \vec{0}
  \]
  \pause
  só é possível com
  \[
    \pause
    \alpha_1 = \cdots = \alpha_k = 0
  \]
\end{frame}

%------------------------------------------------------------------------------%
\begin{frame}{Independência Linear}
  Se
  \[
    \alpha_1\vec{v}_1 + \cdots + \alpha_k\vec{v}_k = \vec{0}
  \]
  \pause
  só é possível com todos \(\alpha_i = 0\)

  \pause
  \bigskip
  Não é possível escrever nenhum vetor como combinação dos demais
\end{frame}

%------------------------------------------------------------------------------%
%------------------------------------------------------------------------------%
\begin{question}
  Determine se os vetores
  \[
    \vec{v}_1 = (1, 2)
  \]
  \[
    \vec{v}_2 = (2, 4)
  \]
  são linearmente dependentes
\end{question}

%------------------------------------------------------------------------------%
\begin{resolution}{Solução}
\begin{columns}[onlytextwidth]
\column{0.5\textwidth}
  \begin{gather*}
    \onslide<+->{\alpha\vec{v}_1+\beta\vec{v}_2 = \vec{0} }
  \\
    \onslide<+->{\alpha\vetor{1}{2} + \beta\vetor{2}{4} = \vetor{0}{0}}
  \\
    \onslide<+->{\begin{cases}
      \alpha+2\beta  = 0 \\
      2\alpha+4\beta = 0
    \end{cases}}
  \\
    \onslide<+->{\alpha = -2\beta}
  \end{gather*}
  \onslide<+->{Escolhendo
  \(
    \beta = 1
  \)}
  \onslide<+->{temos
  \(
    \alpha = -2
  \)}
\column{0.5\textwidth}
  \onslide<+->{Verificando}
  \begin{align*}
    \onslide<+->{\alpha\vec{v}_1+\beta\vec{v}_2}
    \onslide<+->{& = -2\vetor{1}{2} + 1\vetor{2}{4} \\}
    \onslide<+->{& = \vetor{-2}{-4} + \vetor{2}{4} \\}
    \onslide<+->{& = \vetor{0}{0}}
  \end{align*}
  \onslide<+->{Os vetores são linearmente dependentes}
\end{columns}
\end{resolution}

%------------------------------------------------------------------------------%
\endinput

%------------------------------------------------------------------------------%
\begin{question}
  Determine se os vetores
  \[
    \vec{v}_1 = (1, 2)
  \]
  \[
    \vec{v}_2 = (2, 3)
  \]
  são linearmente independentes
\end{question}

%------------------------------------------------------------------------------%
\begin{resolution}{Solução}
\begin{columns}[onlytextwidth]
\column{0.5\textwidth}
  \begin{gather*}
    \onslide<+->{\alpha\vec{v}_1+\beta\vec{v}_2 = \vec{0}}
    \\
    \onslide<+->{\alpha\vetor{1}{2} + \beta\vetor{2}{3} = \vetor{0}{0}}
    \\
    \onslide<+->{\begin{cases}
       \alpha + 2\beta = 0 \\
      2\alpha + 3\beta = 0
    \end{cases}}
  \end{gather*}
  \onslide<+->{Da primeira equação
  \[
    \alpha = -2\beta
  \]}
\column{0.6\textwidth}
  \onslide<+->{Da segunda equação}
  \begin{align*}
    \onslide<+->{2\alpha + 3\beta  & = 0 \\}
    \onslide<+->{2(-2\beta)+3\beta & = 0 \\}
    \onslide<+->{-4\beta+3\beta    & = 0 \\}
    \onslide<+->{-\beta            & = 0 \\}
    \onslide<+->{\beta             & = 0}
  \end{align*}
  \onslide<+->{A única solução é
  \(
    \alpha = \beta = 0
  \)}

  \onslide<+->{Os vetores são linearmente independentes}
\end{columns}
\end{resolution}

%------------------------------------------------------------------------------%
\endinput


%------------------------------------------------------------------------------%
\section{Matrizes e Sistema Lineares}

%------------------------------------------------------------------------------%
\begin{frame}{Sistema Homogêneo}
  Se
  \[
    \pause
    \alpha_1\vec{v}_1 + \cdots + \alpha_k\vec{v}_k = \vec{0}
  \]
  \pause
  Colocando os vetores como colunas de uma matriz
  \[
    \pause
    A =
    \begin{bmatrix}
      \vec{v}_1 & \vec{v}_2 & \cdots & \vec{v}_k
    \end{bmatrix}
  \]
  \pause
  temos
  \[
    \pause
    A
    \begin{pmatrix}
      \alpha_1 \\
      \alpha_2 \\
      \vdots   \\
      \alpha_k
    \end{pmatrix}
    = \vec{0}
  \]
\end{frame}

%------------------------------------------------------------------------------%
\begin{frame}{Critério para Independência Linear}
  Os vetores são linearmente independentes
  \pause
  se o sistema
  \[
    \pause
    A\vec{x} = \vec{0}
  \]
  \pause
  possui somente a solução trivial
  \[
    \pause
    \vec{x} = \vec{0}
  \]
\end{frame}

%------------------------------------------------------------------------------%
%------------------------------------------------------------------------------%
\begin{question}
  Determine se
  \[
    \vec{v}_1 = (1, 0, 2)
  \]

  \[
    \vec{v}_2 = (2, 1, 1)
  \]

  \[
    \vec{v}_3 = (0, 1, 3)
  \]
  são linearmente independentes
\end{question}

%------------------------------------------------------------------------------%
\begin{resolution}{Solução}
\begin{columns}[onlytextwidth]
\column{0.65\textwidth}
  \[\onslide<+->{
    \alpha\vec{v}_1 + \beta\vec{v}_2 + \gamma\vec{v}_3 = \vec{0}}
  \]

  \[\onslide<+->{
      \alpha \vetortri{1}{0}{2}
    + \beta  \vetortri{2}{1}{1}
    + \gamma \vetortri{0}{1}{3}
    = \vetortri{0}{0}{0}}
  \]

  \[\onslide<+->{
    \begin{cases}
      \alpha+2\beta         = 0 \\
      \beta+\gamma          = 0 \\
      2\alpha+\beta+3\gamma = 0
    \end{cases}}
  \]
\column{0.35\textwidth}
    \onslide<+->{Da primeira equação}
    \onslide<+->{\[
      \alpha = -2\beta
    \]}
    \onslide<+->{Da segunda equação}
    \onslide<+->{\[
      \gamma = -\beta
    \]}
\end{columns}
\end{resolution}

%------------------------------------------------------------------------------%
\begin{resolution}{Solução}
\begin{columns}[onlytextwidth]
\column{0.5\textwidth}
  \onslide<+->{Substituindo na terceira}
  \begin{align*}
    \onslide<+->{2\alpha+\beta+3\gamma      & = 0 \\}
    \onslide<+->{2(-2\beta)+\beta+3(-\beta) & = 0 \\}
    \onslide<+->{-4\beta+\beta-3\beta       & = 0 \\}
    \onslide<+->{-6\beta                    & = 0 \\}
    \onslide<+->{\beta                      & = 0}
  \end{align*}
\column{0.5\textwidth}
  \begin{align*}
    \onslide<+->{\alpha & = -2\beta = 0 \\}
    \onslide<+->{\gamma & =  -\beta = 0}
  \end{align*}
  \onslide<+->{Portanto
  \(
    \alpha=\beta=\gamma=0
  \)}

  \onslide<+->{e os vetores são linearente independentes}
\end{columns}
\end{resolution}

%------------------------------------------------------------------------------%
\endinput


%------------------------------------------------------------------------------%
\begin{frame}{Determinante}
  Dados os vetores
  \[
    \pause
    \vec{v}_1, \vec{v}_2, \ldots, \vec{v}_n \in \R^n
  \]
  \pause
  formamos a matriz quadrada
  \[
    \pause
    A =
    \begin{bmatrix}
      \vec{v}_1 & \vec{v}_2 & \cdots & \vec{v}_n
    \end{bmatrix}
  \]
  \pause
  Se
  \[
    \det(A) \neq 0
  \]
  \pause
  os vetores são linearmente independentes
\end{frame}

%------------------------------------------------------------------------------%
\begin{frame}{Vetor nulo}
  Qualquer conjunto que contenha
  \[
    \vec{0}
  \]
  \pause
  é linearmente dependente
\end{frame}

%------------------------------------------------------------------------------%
\begin{frame}{Múltiplos}
  Se
  \[
    \vec{v}_2 = \alpha\vec{v}_1
  \]
  \pause
  para algum
  \[
    \alpha\in\R
  \]
  \pause
  então
  \[
    \vec{v}_1
    \quad\text{e}\quad
    \vec{v}_2
  \]
  \pause
  são linearmente dependentes
\end{frame}

%------------------------------------------------------------------------------%
%------------------------------------------------------------------------------%
\begin{question}
  Determine se
  \[
    \vec{v}_1 = (1, 2, 0)
  \]

  \[
    \vec{v}_2 = (0, 1, 1)
  \]

  \[
    \vec{v}_3 = (2, 0, 1)
  \]
  são linearmente independentes
\end{question}

%------------------------------------------------------------------------------%
\begin{resolution}{Solução}
  \[
    \onslide<+->{A
    \;=\; \begin{pmatrix}
          \vec{v}_1 & \vec{v}_2 & \vec{v}_3
        \end{pmatrix}}
    \onslide<+->{\;=\; \begin{pmatrix}
          1 & 0 & 2 \\
          2 & 1 & 0 \\
          0 & 1 & 1
        \end{pmatrix}}
  \]
  \bigskip
  \[
    \onslide<+->{\det(A)}
    \onslide<+->{= 1
      \begin{vmatrix}
        1 & 0 \\
        1 & 1
      \end{vmatrix}
      - 0
      + 2
      \begin{vmatrix}
        2 & 1 \\
        0 & 1
      \end{vmatrix}}
    \onslide<+->{= 1 + 4}
    \onslide<+->{= 5}
  \]

  \onslide<+->{
  \bigskip
  Como
  \(
    \det(A) = 5 \neq 0
  \)}
  \onslide<+->{os vetores são linearmente independentes}
\end{resolution}

%------------------------------------------------------------------------------%
\endinput

%------------------------------------------------------------------------------%
\begin{question}
  Determine se
  \[
    \vec{v}_1 = (1, 2, 1)
  \]

  \[
    \vec{v}_2 = (2, 1, 0)
  \]

  \[
    \vec{v}_3 = (3, 3, 1)
  \]
  são linearmente independentes
\end{question}

%------------------------------------------------------------------------------%
\begin{resolution}{Solução}
\begin{columns}
\column{0.5\textwidth}
  \[
    \onslide<+->{\alpha\vec{v}_1 + \beta\vec{v}_2 + \gamma\vec{v}_3 = \vec{0}}
  \]

  \[
    \onslide<+->{\begin{cases}
      \alpha+2\beta+3\gamma = 0 \\
      2\alpha+\beta+3\gamma = 0 \\
      \alpha+\gamma         = 0
    \end{cases}}
  \]
  \onslide<+->{Da terceira equação
  \[
    \alpha = -\gamma
  \]}
\column{0.5\textwidth}
  \onslide<+->{Substituindo na primeira}
  \begin{align*}
    \onslide<+->{\alpha+2\beta+3\gamma  & = 0       \\}
    \onslide<+->{-\gamma+2\beta+3\gamma & = 0       \\}
    \onslide<+->{2\beta+2\gamma         & = 0       \\}
    \onslide<+->{\beta                  & = -\gamma}
  \end{align*}
\end{columns}
\end{resolution}

%------------------------------------------------------------------------------%
\begin{resolution}{Solução}
\begin{columns}
\column{0.5\textwidth}
  \onslide<+->{Substituindo na segunda}
  \begin{align*}
    \onslide<+->{2\alpha+\beta+3\gamma        & = 0  \\}
    \onslide<+->{2(-\gamma)+(-\gamma)+3\gamma & =  0 \\}
    \onslide<+->{0                            & = 0}
  \end{align*}
  \onslide<+->{Logo $\gamma$ é livre}

  \onslide<+->{Tomando
  \(
    \gamma = 1
  \)}
  \onslide<+->{obtemos}
  \[
    \onslide<+->{\alpha = -\gamma = -1}
    \qquad
    \onslide<+->{\beta = -\gamma = -1}
  \]
\column{0.5\textwidth}
  \onslide<+->{Temos a relação}
  \[
    \onslide<+->{-\vec{v}_1-\vec{v}_2+\vec{v}_3 = \vec{0}}
  \]
  \[
    \onslide<+->{\vec{v}_3 = \vec{v}_1 + \vec{v}_2}
  \]
  \onslide<+->{Portanto os vetores são linearmente dependentes}
\end{columns}
\end{resolution}

%------------------------------------------------------------------------------%
\endinput


%%------------------------------------------------------------------------------%
%\section{Lista Mínima}
%
%%------------------------------------------------------------------------------%
%\begin{frame}{Lista Mínima}
%  \begin{enumerate}
%    \item
%  \end{enumerate}
%
%  \vfill
%  Atenção: A prova é baseada no livro, não nas apresentações
%\end{frame}

%------------------------------------------------------------------------------%
\end{document}
%------------------------------------------------------------------------------%
